\documentclass{article}
\usepackage{decomp}
\begin{document}
\sourcefile{elife\_figure\_supplements/ (fig1-figsupp1.jpg … fig2-figsupp6.jpg, elife\_api\_article.json)}
\decompnote{Figure supplements of Templier T (2019) eLife 8:e45696. They are not in the article PDF, which prints only their titles and DOIs. Images: the full-size JPEGs from iiif.elifesciences.org (copied unchanged to \texttt{figures/sfig1-1.jpg} … \texttt{figures/sfig2-6.jpg}). Legends: verbatim from the eLife API record of the article (\texttt{elife\_api\_article.json}, objects \texttt{fig1s1} … \texttt{fig2s6}: label, title, caption), with only the HTML markup turned into LaTeX (\texttt{work/legends/make\_legends.py}; checked back against the JSON by \texttt{work/legends/check\_legends.py}). There are no PDF pages, so there are no \texttt{\textbackslash pdfpage} marks. The eLife web page gives each supplement after its main figure; that order is kept. The decompiler's file names: Figure N---figure supplement M = \texttt{sfigN-M}, panel x = \texttt{sfigN-Mx}.}

\decompnote{Figure supplements of Figure 1.}

\begin{center}
\includegraphics[width=0.75\linewidth,height=0.5\textheight,keepaspectratio]{figures/sfig1-1.jpg}
\end{center}
\noindent\textbf{Figure 1—figure supplement 1.} \textbf{Wafer overview for~Dataset 2.} (\textbf{a}) Widefield LM of the collected sections. The intensity differences come from section thickness inhomogeneity. As documented by Harris et al. (2006)~, the lack of an hermetic enclosure during sectioning typically induces section inhomogeneity. (\textbf{b}) Close-up around one section, showing~merged multichannel fluorescent imagery: blue~for~Coumarin, with green- and red-fluorescent beads. (\textbf{c})~Montage of all sections with the same orientation. Scale bars: (a)~4 mm; (b)~500 μm; (c)~2 mm.

\noindent DOI: https://doi.org/10.7554/eLife.45696.003
\alttext{Three light-microscopy panels. (a) Greyscale widefield overview of the wafer: many small rectangular sections (Dataset 2 has 203 per the main text), darker and lighter grey (each with a lighter and a darker half), scattered in a broad ring with a sparser centre, on a light background with a faint regular tile pattern from the mosaic; black margins on the right and at the edges; black scale bar at lower right (4 mm per the legend). (b) Close-up of one section in the centre and parts of its neighbours: each section has a blue half densely speckled with small cyan/green and a few red dots (the magnetic resin with beads) and a dark grey half (tissue); white background; black scale bar at lower right (500 μm). (c) Montage of all sections in a regular grid of about 12 columns and 17 rows, each tile a section in the same orientation with the blue (coumarin) part at the top and the grey part below; the last row is incomplete; black scale bar at lower right (2 mm).}

\begin{center}
\includegraphics[width=0.85\linewidth,height=0.5\textheight,keepaspectratio]{figures/sfig1-2.jpg}
\end{center}
\noindent\textbf{Figure 1—figure supplement 2.} \textbf{Electron micrographs of the sections for~Dataset 1.} (\textbf{a}) Electron micrograph of numerous sections collected on wafer. Most of the collected sections exhibit a dummy (like that highlighted in the yellow box) that is attached to their tissue part and that comes from the previous section, which is explained by the following mechanism: with this sample block, the dummy part often detached from the section during sectioning or immediately after landing on the water surface. The main part (magnetic resin and tissue) of the last cut section was then floating freely while its dummy remained attached to the edge. Then the main part of the next cut section stuck to the previous dummy and so on. (\textbf{b}) EM of a section. The yellow dashed square highlights the region that has been imaged with the electron microscope and became darker due to the beam irradiation. (\textbf{c}) EM of well-dispersed superparamagnetic nanoparticles in the appended resin. (\textbf{d}) Small contaminations with the~appearance of oily droplet can sometimes be found in the magnetic resin. These could be unreacted oleic acids,~an~oily compound from handling during preparation, or something unknown to the author. Scale bars: (a)~1~mm; (b)~100 µm; (c)~500 nm; (d)~2~µm.

\noindent DOI: https://doi.org/10.7554/eLife.45696.004
\alttext{Four greyscale electron micrographs. (a) Low magnification: dozens of sections on a black background (the wafer), each a rectangle with a light grey uniform part and a darker textured part (tissue, with white blood-vessel profiles); many sections carry a thin light strip at one edge; one such strip is outlined by a small yellow box just left of centre; white scale bar at lower right (1 mm). (b) One section: a uniform light upper part (magnetic resin) and a darker lower part with grey round profiles and white vessel profiles (tissue); a yellow dashed square in the tissue marks a region that is slightly darker than its surroundings; a light strip along the left and bottom edges; black scale bar at lower right (100 µm). (c) Light grey field densely sprinkled with small dark dots of slightly different size and darkness, evenly spread (nanoparticles); black scale bar at lower right (500 nm). (d) Light grey grainy field with one round dark-rimmed droplet with a grey interior near the top centre and a few small dark specks; black scale bar at lower right (2 µm).}

\begin{center}
\includegraphics[width=0.7\linewidth,height=0.55\textheight,keepaspectratio]{figures/sfig1-3.jpg}
\end{center}
\noindent\textbf{Figure 1—figure supplement 3.} \textbf{Magnetic augmentation.} (\textbf{a})~Mounting helper device. The manual manipulator allows the experimenter to precisely~place the needle in contact with the block to be glued onto the sample. The iron plate together with the base magnet of the needle holder maintain the position set by the manipulator. (\textbf{b}) The detachable part is placed in an oven for temperature curing. (\textbf{c}) The biological sample is trimmed manually with a razor blade. (\textbf{d}) The magnetic resin is glued to the sample and maintained in place with a needle. (\textbf{e}) Overlay of color and fluorescent imagery showing the fluorescent particles contained in~the magnetic resin. (\textbf{f}) The magnetic resin is manually trimmed down to achieve roughly a 50/50 ratio of sample surface to magnetic surface suitable for sectioning at 50 nm nominal thickness. (\textbf{g}) An additional dummy piece of heavy-metal-stained brain tissue embedded in resin is glued to the block. Scale bars: (\textbf{a},\textbf{~b})~20 mm; (\textbf{c–g}): 500 μm.

\noindent DOI: https://doi.org/10.7554/eLife.45696.005
\alttext{Seven colour photographs. (a) The mounting helper device on an aluminium base, side view, with white labels and arrows: ``Sample'' (a small resin block on a holder at the far left), ``Needle'' (a thin needle touching it), ``Link'' with the subtitle ``Needle holder - Manipulator'' (a horizontal rod), ``Manual manipulator'' (a black Narishige micromanipulator with large knobs, right), ``Magnet'' (a black disc base of the needle holder), ``Iron plate'', and ``Detachable part'' (the block below the plate, outlined by a white dashed line); small arrows marked x and y show two directions on the plate; black scale bar at lower right (20 mm per the legend). (b) The detachable part with the needle holder and sample standing on the wire shelf of an oven; white scale bar (20 mm). (c) Top view of a trimmed sample block face: a light, cyan-green square with fine tissue texture inside a dark frame on a green background; white scale bar (500 μm). (d) The dark sample block (left) with a larger dark brown block glued to its right side (magnetic resin), held by a needle coming in from the right; black scale bar (500 μm). (e) The same view in a colour--fluorescence overlay: the right block is densely covered with green and red dots; black scale bar (500 μm). (f) The magnetic block trimmed down to a narrower rectangle next to the sample block; black scale bar (500 μm). (g) An additional dark piece glued at the right of the brown magnetic resin, held by a needle from the right; black scale bar (500 μm).}

\begin{center}
\includegraphics[width=0.75\linewidth,height=0.5\textheight,keepaspectratio]{figures/sfig1-4.jpg}
\end{center}
\noindent\textbf{Figure 1—figure supplement 4.} \textbf{Wafer overview of the 100 sections collected during the video-recorded session.} Scale bar: 5 mm.

\noindent DOI: https://doi.org/10.7554/eLife.45696.006
\alttext{Greyscale widefield light-microscopy overview (a mosaic with visible tile shading) of the small rectangular sections (100 per the title) on the wafer, each with a lighter and a darker part, lying in a broad U-shaped band from the upper left down across the middle and up to the right, the centre and the top right being free of sections; a few thin dark strips (debris or section fragments) among them; a small faint ring at the lower left; black scale bar at lower right (5 mm per the legend). No plotted data.}

\begin{center}
\includegraphics[width=0.95\linewidth,height=0.5\textheight,keepaspectratio]{figures/sfig1-5.jpg}
\end{center}
\noindent\textbf{Figure 1—figure supplement 5.} \textbf{Multibeam scanning EM of magnetically collected sections on silicon wafer.} (\textbf{a}) Overview of the 15 consecutive sections collected with MagC that were used for the first MultiSEM experiment.~(\textbf{b}) Overview of 91 stitched tiles. (\textbf{c}) A tile produced by one of the 91 beams. The image quality is representative of the quality observed in the whole section. (\textbf{d}) An inset from a tile acquired with a Zeiss MultiSEM of a section from Dataset 1 after the complete CLEM pipeline and an additional wafer-wide broad ion beam milling of about 25 nm (Templier, in preparation). Imaging conditions: silicon wafer chip glued to the~EM stub with carbon glue, 4 nm pixel size, 400 ns dwell time, Scale bars: (\textbf{a}) 2 mm; (\textbf{b}) 20 μm; \textbf{c}) 2 μm; \textbf{d}) 1 μm.

\noindent DOI: https://doi.org/10.7554/eLife.45696.007
\alttext{Four panels left to right, linked by zoom lines. (a) Greyscale light-microscopy overview of 15 sections lying in two slanted, partly overlapping columns on the wafer chip; a small green square on one section is joined by green lines to panel b; black scale bar (2 mm per the legend). (b) Greyscale EM overview of the stitched multibeam field: a roughly hexagonal mosaic of square tiles with stepped edges on a black background, green frame; a small yellow square near the centre is joined by yellow lines to panel c; white scale bar (20 μm). (c) EM tile with neuropil: densely packed membrane profiles, myelinated axons at the top, mitochondria; yellow frame; black scale bar (2 μm). (d) Higher-magnification EM inset of neuropil with clearly outlined membranes, mitochondria with cristae and synaptic vesicles; black scale bar (1 μm). No plotted data.}

\decompnote{Figure supplements of Figure 2.}

\begin{center}
\includegraphics[width=0.75\linewidth,height=0.5\textheight,keepaspectratio]{figures/sfig2-1.jpg}
\end{center}
\noindent\textbf{Figure 2—figure supplement 1.} \textbf{Section order retrieval for~Dataset 1 (low concentration of beads).} (\textbf{a}) Matrix of pairwise similarities of unordered sections computed with EM imagery. Darker pixels depict higher similarity whereas~white pixels depict no similarity. (\textbf{b}) Reordered EM matrix. (\textbf{c}) Matrix of pairwise similarities of unordered sections computed with fluorescent bead imagery. The original order is the order provided by the section segmentation pipeline. (\textbf{d}) Similarity matrix of the reordered sections. The order overall looks to~be~consistent, except for~slight deviations at the end of the dataset (around section number 500) that can be seen in the lower left of the matrix.

\noindent DOI: https://doi.org/10.7554/eLife.45696.009
\decompnote{Legend and image do not match for panels c and d. The image shows: (a) an unordered similarity matrix, (b) a reordered matrix (dark diagonal band), (c) a reordered matrix with a dark diagonal band and many thin coloured horizontal and vertical lines, and (d) a bar chart with values printed over the bars, with a colour legend (1, 2, 3, 4, 7) between panels b and d. The legend describes (c) as the unordered bead matrix and (d) as the reordered matrix. The legend's ``lower left of the matrix'' for deviations around section 500: section 500 is at the lower right end of the diagonal.}
\alttext{Four panels in a 2 × 2 layout. Panels a, b and c are square similarity matrices drawn as greyscale images in black frames; top and left axes with ticks every 100 and labels 0, 200, 400 (top labels on a and b, left labels on a and c; no axis titles). The frame ends a little after the 500 tick: the matrices span about 514 sections (measured on the image; Dataset 1 has 507 sections per the main text). The panel letters are printed in large black type at the upper right inside each matrix. (a) Scattered small grey and black dots over the whole white field and a thin black main diagonal; the dots are denser in places near the bottom. (b) White except for a dark band along the main diagonal, black at the centre and grey a few tens of sections to either side; at the lower right end (around section 500) the band widens and gets lighter, with a small irregular patch. (c) A dark diagonal band like in b, overlaid by many thin semi-transparent straight lines running across the whole matrix, horizontally and vertically at the same positions: mostly green, many blue, and single magenta (at about section 250) and red (at about section 488, near the right/bottom edge) lines; some lines are close together in groups. Between b and d a colour legend: green line ``1'', blue ``2'', dark grey ``3'', magenta ``4'', red ``7''; below it a horizontal black rule. (d) Bar chart, no axis titles: x categories 0, 1, 2, 3, 4, 5, 6, 7; y axis 0 to 400 with labels every 100; blue bars with the value printed above each: 375, 103, 26, 7, 1, 0, 0, 1 (the bars for 4 and 7 are thin slivers, those for 5 and 6 absent).}
\figuredata{a}{figures/sfig2-1a.csv}{figures/sfig2-1.py}{digitized from the eLife JPEG (tier 3): grey level per image pixel (one pixel ≈ 1.35 matrix cells, cells not resolved), section index of each pixel from the tick marks; uncalibrated (no colour bar)}
\figuredata{b}{figures/sfig2-1b.csv}{figures/sfig2-1.py}{as panel a (tier 3)}
\figuredata{c}{figures/sfig2-1c.csv \textup{and} figures/sfig2-1c\_lines.csv}{figures/sfig2-1.py}{as panel a (tier 3); the coloured lines are listed separately with position and legend colour, classified from the JPEG colours (lower bound: adjacent lines merge; ±1.5 sections)}
\figuredata{d}{figures/sfig2-1d.csv}{figures/sfig2-1.py}{bar values printed in the figure (exact, transcribed) plus bar heights digitized from the JPEG (tier 3, ±1.7)}

\begin{center}
\includegraphics[width=0.6\linewidth,height=0.55\textheight,keepaspectratio]{figures/sfig2-2.jpg}
\end{center}
\noindent\textbf{Figure 2—figure supplement 2.} \textbf{Automated LM-EM registration.} (\textbf{a}) Histogram of the~numbers of matching inliers found for each of the 203 LM-EM pairs of Dataset 2. (\textbf{b}) A reflection brightfield light micrograph after simple thresholding. (\textbf{c})~Downscaled EM mosaic. (\textbf{d}) Same micrograph as in (a) after local contrast normalization. Note the high similarity with its EM counterpart micrograph in (c). (\textbf{e,~f}) Same micrographs as in (c)~and~(d), respectively. The yellow crosses show the locations of~SIFT features that match between the two images. The dashed yellow lines in (e) show the outline of the~LM micrograph when affine transformed to match its EM counterpart. Scale bars: 50 μm.

\noindent DOI: https://doi.org/10.7554/eLife.45696.010
\decompnote{Panel (d) is described as ``Same micrograph as in (a)''; panel (a) is the histogram, the micrograph is in (b).}
\alttext{Six panels in three rows of two. (a) Histogram: x axis ``Number of inliers'', linear, ticks every 500 with labels 0 (under the y axis), 1000, 2000; y axis ``Frequency'', ticks every 2 with labels 0, 4, 8; 60 grey bars with black outlines from about 67 to about 2545 inliers (bins about 41 wide). Heights range from 0 to 10; the distribution is broad, with most bars between 2 and 7, a maximum of 10 in the bin at about 1595--1636 inliers and a second peak of 8 at about 1430--1471; empty bins at about 108--149, 397--439, 1802--1843, 2173--2215 and 2421--2504 (values read off the image). The 60 bar heights read off the image sum to 201. (b) Reflection brightfield light micrograph (greyscale), labelled ``LM'' at lower right: a mottled grey field with lighter round and elongated profiles and a large light ring-shaped profile (a blood vessel) in the centre; uneven illumination, darker towards the bottom; white scale bar at lower left (50 μm). (c) Downscaled EM mosaic labelled ``EM'': the same field in EM contrast with many light round profiles (cell bodies and vessels) and the large central vessel with a dark lumen; black frame edges; black scale bar at lower left. (d) The LM micrograph after local contrast normalization, labelled ``LM'': it now looks much like the EM image; black scale bar at lower left. (e) The EM mosaic in cyan, labelled ``EM'', with hundreds of small yellow crosses spread over it and a yellow dashed outline (a tilted rectangle) marking the transformed LM field. (f) The normalized LM image in cyan, labelled ``LM'', with the yellow crosses at the matching features.}
\figuredata{a}{figures/sfig2-2a.csv}{figures/sfig2-2a.py}{digitized from the eLife JPEG (tier 3): 60 bins, bar heights read within 0.02 of integers and rounded; x from the 500-step ticks}

\begin{center}
\includegraphics[width=0.6\linewidth,height=0.6\textheight,keepaspectratio]{figures/sfig2-3.jpg}
\end{center}
\noindent\textbf{Figure 2—figure supplement 3.} \textbf{Sectioning quality and restart of a MagC block.All images are SEM micrographs of sections manually collected on a silicon wafer.} (\textbf{a}) Section number \#100. The yellow box indicates the region acquired at high resolution for the sections up to number \#\textasciitilde{}3500 (see panel (\textbf{c}) for sections above \#4000). (\textbf{b}) Montage of the high-resolution regions acquired. The corresponding indices are shown on the right. The insets numbered 1 to 4 are shown in panel (\textbf{d}). Insets 1, 2 and 3 are the only locations where vertical streaks resulting from sectioning were found. Inset 4 in section \#\textasciitilde{}5000 is representative of the streak-free sectioning quality manually observed in all acquired high-resolution regions. The insets illustrating restart continuity (labeled \textbf{e}, \textbf{f}, and \textbf{g}) are shown below in panels with corresponding lettering. (\textbf{c}) Complete or partial sections collected after the restart. The yellow box shows the region acquired in section numbers above \#4000. The nature of the cuts was observed through the binocular microscope simultaneously during the sectioning restart and is as follows (\#1 refers to \#\textasciitilde{}4000 + 1 and so on): \#1, partial tissue; \#2, partial tissue; \#3, partial tissue; \#4, partial tissue connected with partial magnet; \#5, partial tissue connected with partial magnet (but disconnected during manual pickup, the magnetic and dummy parts remained stuck to the bottom of section \#6); \#6, partial tissue with disconnected partial magnet (the tissue/ magnet separation occurred immediately upon sectioning); \#7, full tissue connected with full magnet (but disconnected during manual pickup, the magnetic portion is visible immediately to the right of the section); \#8 to \#15, full normal sections (not collected). (\textbf{d}) Insets 1, 2 and 3 from panel (\textbf{b}) showing (with yellow arrows) the only vertical streaks found after manual inspection among the entire high-resolution imagery from panel (\textbf{b}). Inset 4 from section \#\textasciitilde{}5000 shows a representative image of the lack of image quality impairment that might occur due to a potentially damaged knife. (\textbf{e– g}) Four consecutive insets from the restart sections shown in yellow in panel (\textbf{b}). The interior of three small neuronal processes are painted in red, green and blue. There is no issue in tracing fine processes across the 20 hr long restart. The good continuity in panels (\textbf{e–g}) indicates that no material was lost despite the partial sections. Scale bars: (\textbf{a}) 200 mm; (\textbf{b}) 100 mm; (\textbf{c}) 500 mm; (\textbf{d–g}) 2 mm.

\noindent DOI: https://doi.org/10.7554/eLife.45696.011
\decompnote{The scale bars are given in ``mm'' (200, 100, 500 and 2 mm). The micrographs show single sections and ultrastructure, so μm is almost certainly meant: the μ of the source seems to have become ``m'' (as in the PDF's own text layer). Kept as in the JSON.}
\alttext{SEM micrographs, panels a--g. (a) A whole section labelled ``\# 100'' at the top right: uniform light upper part, textured tissue lower part with white vessel profiles, a light strip at the bottom; a thin yellow horizontal box across the tissue part; white scale bar at lower right. (b) A stack of long horizontal image strips (the high-resolution bands), each labelled on the right: ``\#\textasciitilde100'', ``\#\textasciitilde500'', ``\#\textasciitilde1000'', ``\#\textasciitilde1500'', ``\#\textasciitilde2000'', ``\#\textasciitilde2500'', ``\#\textasciitilde3000'', ``\#\textasciitilde3500'', ``\#\textasciitilde4000'', then ``\#(\textasciitilde4000) + 1'' to ``\#(\textasciitilde4000) + 7'' and ``\#\textasciitilde 5000''. Arrows and notes on the far right: ``Knife cleaning'' pointing between the \#\textasciitilde2000 and \#\textasciitilde2500 strips; ``Knife cleaning / Sham collection / 20h break / Restart'' pointing between \#\textasciitilde4000 and \#(\textasciitilde4000) + 1; ``Visual inspection: no section miss'' with a bracket between \#(\textasciitilde4000) + 7 and \#\textasciitilde5000. The restart strips +1 to +5 are partial (shorter, on black), +6 and +7 full. Small yellow squares numbered 1 (strip \#\textasciitilde1000), 2 and 3 (strip \#\textasciitilde2000) and 4 (strip \#\textasciitilde5000) mark insets; three vertical yellow lines with small squares, labelled e, f and g on the \#\textasciitilde4000 strip, run down through the restart strips (e to +3, f to +5, g to +7). White scale bar on the left. (c) Seven SEM images of the restart cuts side by side, labelled below ``\#(\textasciitilde4000) + 1'' to ``\#(\textasciitilde4000) + 7'': +1 and +2 small irregular tissue pieces, tilted, on black; +3 to +7 larger section parts with a grey tissue rectangle and lighter magnetic parts, with pieces attached or separate in +4 and +6; a thin yellow horizontal box runs across all seven at the same height; white scale bar at upper right. (d) Four EM insets numbered 1--4 at upper right: in 1 and 2 a pair of yellow arrowheads (top and bottom edge) marks a faint vertical streak; in 3 two pairs mark two closely spaced streaks; 4 shows clean neuropil with no streak; black scale bar in inset 4. (e), (f), (g) Three rows of four consecutive EM insets each, showing the same neuropil across the restart; in every inset three small profiles are filled red, green and yellow at matching positions (the legend says red, green and blue); (g) contains a large dark myelinated axon in the centre; white scale bar at lower right of each row.}

\begin{center}
\includegraphics[width=0.75\linewidth,height=0.6\textheight,keepaspectratio]{figures/sfig2-4.jpg}
\end{center}
\noindent\textbf{Figure 2—figure supplement 4.} \textbf{Labeled axon making a synapse en passant.} The axon is delineated with dashed yellow lines (every second section). Black arrows indicate the synapse. Scale bar:~1~µm.

\noindent DOI: https://doi.org/10.7554/eLife.45696.012
\alttext{A gallery of 48 consecutive EM images in 8 rows of 6 (read left to right, top to bottom), all the same field of neuropil, greyscale, with the green fluorescence of the tracer overlaid as a soft green glow. In the upper rows a large dark myelinated axon profile sits at the upper right. A small profile at the centre-left grows into a larger round profile (a bouton) filled with green signal, and then shrinks again in the lower rows; in many tiles it is outlined with a dotted yellow line. Black arrows pointing to the upper right mark the synapse at the left edge of the bouton in 11 tiles (the last three tiles of row 2, all six of row 3 and the first two of row 4). Black scale bar in the last tile (1 µm per the legend).}

\begin{center}
\includegraphics[width=0.7\linewidth,height=0.6\textheight,keepaspectratio]{figures/sfig2-5.jpg}
\end{center}
\noindent\textbf{Figure 2—figure supplement 5.} \textbf{Section collection quality.} (\textbf{a}) Schematic of the immersed tilted silicon wafer chip. Side view from the experimenter looking towards the ultramicrotome (same point of view as in Figure 1b). A few microscopy coverslips are asymmetrically stacked to slightly tilt the wafer chip in~order~to avoid accumulation of water surface dust in the center of the wafer at the end of the water withdrawal. (\textbf{b}) The single tear found in Dataset 2, which was present in the magnetic portion of section \#22. (\textbf{c–e}) tear in Dataset one in sections \#258, \#359 and \#481, respectively. (\textbf{f})~Schematic of the interaction of water with a hydrophobic substrate, a partially hydrophobic substrate, and a substrate with nearby walls. All three configurations present a~risk of impairing sections upon drying.

\noindent DOI: https://doi.org/10.7554/eLife.45696.013
\alttext{Six panels. (a) Schematic, side view: a dashed line labelled ``Water level'' at the top; below it a grey bar labelled ``Tilted silicon wafer chip'', tilted up to the right, with a short black line and an arrow labelled ``\textasciitilde2 degrees'' marking the angle at its right end; the bar rests on two stacks of thin black lines (two lines on the left, five on the right) labelled ``Asymmetrically stacked coverslips'', which stand on a cyan block labelled ``Knife boat''. (b) Fluorescence image (black background) of the magnetic portion of one section, densely dotted with red and green beads; a smooth green, bead-free tongue intrudes from the bottom centre and a lighter band runs along the lower right edge (the tear region); white-boxed black scale bar ``100 μm'' at lower right. (c), (d), (e) Greyscale EM images of tissue (cell bodies, vessels) each crossed top to bottom by a straight narrow grey band (the tear) without tissue texture, near the left in c and e and in the centre in d; scale bar ``50 μm'' at lower right of e. (f) Schematic: a grey bar labelled ``Wafer'' with a dashed black surface line; three cyan water profiles on it: a high rounded drop at the left, a lower asymmetric wave in the middle, and at the right a surface rising steeply against a black vertical wall.}

\begin{center}
\includegraphics[width=0.85\linewidth,height=0.5\textheight,keepaspectratio]{figures/sfig2-6.jpg}
\end{center}
\noindent\textbf{Figure 2—figure supplement 6.} \textbf{Multicolor correlative LM-EM imagery (Dataset 1) visualized in neuroglancer.} Multicolor correlative LM-EM imagery (Dataset 1) of zebra finch HVC nucleus showing three neuroanatomical tracers injected in Area X (blue), the nucleus Robustus of the Arcopallium (green), and Avalanche (red). The two panes on the right show x and y re-slices through the volume.

\noindent DOI: https://doi.org/10.7554/eLife.45696.014
\alttext{Screenshot of the neuroglancer viewer, four panes. Upper left: an EM section of zebra finch HVC in greyscale with blue, red and green fluorescent signal overlaid (blue blobs and a large blue patch near the centre, red spots around two cell bodies, green rings at the lower left); a white scale bar ``15 μm'' at lower left; thin red and green cursor lines. Upper right: a re-slice through the volume shown as a horizontal strip on a grey background, with a large cell body outlined by blue signal and a red and blue cursor. Lower left: a 3D view (``Slices'' checkbox at upper left) showing three intersecting orthogonal planes of the volume on black. Lower right: the other re-slice as a vertical strip on grey with green and blue cursor lines. ``POWERED BY NEUROGLANCER'' at the lower right corner.}

\end{document}
